\section{Cálculo Numérico}

\subsection{Errores Computacionales}

Si usted ha intentado realizar una suma de dos números decimales mediante algún lenguaje de programación, podrá haber notado que el resultado no es exactamente correcto
\begin{sourcecode}{python}{Suma de dos decimales en Python}
  a = 0.1
  b = 0.2
  c = 0.3

  print(f"{a + b == c}")
  print(f"{a + b}")
\end{sourcecode}
Cuyo output será

\begin{sourcecode}{plaintext}{Output}
  False
  0.30000000000000004
\end{sourcecode}

Este error se debe a la manera en que las computadoras deben representar números. En general, los humanos representamos los números en base 10, por ejemplo, escribimos el 325 como
\begin{equation*}
  325 = 3\cdot 10^{2} + 2\cdot 10^{1} + 5\cdot 10^{0} = [325]_{10}
\end{equation*}
Pero los computadores utilizan bits (0 y 1), por lo que estos representan números en base 2, es decir:
\begin{equation*}
  325 = 2^{8} + 2^{6} + 2^{2} + 2^{0} = [101000101]_{2}
\end{equation*}

Ahora, para representar número reales en el sistema binario, se deben incluir exponentes negativos a la suma de exponentes. En la representación binaria, los números se separan por un exponente y una parte fraccionaria (mantisa) las cuales son separadas por un punto ($.$). Por ejemplo,

\begin{align*}
  [2.145]_{10} &= 2\cdot 10^{0} + \frac{1}{8} \\
               &= 2^{1} + 2^{-3} \\
               &= [10.001]_{2} \\
               &= [1.0001]_{2} \cdot 2^{1}
\end{align*}

Este ejemplo en particular funciona bien porque la mantisa del númer decimal resulta ser un exponente de 2, veamos que ocurre para un número más complejo:

\begin{align*}
  [0.1]_{10} &= [1]_{2}/[1010]_{2} \\
             &= [0.0\overline{0011}]_{2} \\
             &= [0.00011001100110011 \cdots] \\
             &= 2^{-4} + 2^{-5} + 2^{-8} + 2^{-9} + 2^{-12} + 2^{-13} + \cdots \\
             &= 2^{-4}\cdot (1+ 2^{-1} + 2^{-4} + 2^{-5} + \cdots)
\end{align*}
Obtenemos una serie infinita de potencias de dos. Como en la practica los computadores tienen memoria finita es necesario truncar la serie. En general, la mantisa recibe 52 bits, por lo que la serie se trunca hasta el termino $2^{-52}$. en efecto:
\begin{align*}
  [0.1]_{10} &= 2^{-4}\cdot(1+ 2^{-1} + 2^{-4} + 2^{-5} + \cdots + 2^{-52})
\end{align*}

Este número no será \textit{exactamente} igual a $0.1$, si no que será igual a $0.1 \cdot (1+ \epsilon)$. Este será el caso para la gran mayoría de los calculos hechos computacionalmente, el resultado final tendrá una pequeña divergencia con respecto al valor real, es decir 
$$\bar{a} = a\cdot (1+ \epsilon) $$
donde $\bar{a}$ y $a$ son los valores calculado y real, respectivamente. La diferencia entre el número deseado y el resultante se denomina \textbf{error computacional}. Dado que las computadores no pueden físicamente lidiar con cantidades infinitesimalmente pequeñas, muchos de los métodos de cálculo adaptados a la programación que veremos a lo largo del apunte serán basados en aproximaciones, por esta razón es importante aprender a identificar fuentes de errores numericos y sus ordenes de magnitud.

\subsection{Derivadas numéricas}

Sabemos del calculo diferencial que la derivada de una función arbitraria $f$ evaluada en un punto arbitrario $x_{0}$  es dada por
\begin{equation}
  f^{\prime}(x_{0}) = \lim_{h \to 0} \frac{f(x_{0}+h) - f(x_{0})}{h}
\end{equation}
pero ya sabemos que necesitamos que $h$ sea un número finito, por lo que estamos obligados a seguir anaizando esta definición. De la serie de taylor de $f$ evaluada en un punto $a + h$, $h > 0$, tenemos que:

\begin{equation}
  f(a + h) = f(a) + f^{\prime}(a)h + \frac{1}{2!}f^{\prime \prime}(a)h^{2} + \frac{1}{3!}f^{\prime \prime \prime}(a)h^{3} + \cdots
\end{equation}
Podemos reemplazar los terminos de orden superior por un $f(\xi)h^{2}$, $a \leq \xi \leq a + h$ (Teorema del valor medio), luego
\begin{equation}
  f(a + h) = f(a) + f^{\prime}(a)h + \frac{1}{2!}f^{\prime \prime}(\xi)h^{2}
\end{equation}
\begin{equation} \label{eq:derad}
  f^{\prime}(a) = \frac{f(a + h) - f(a)}{h} - \frac{1}{2!}f^{\prime \prime}(\xi)h
\end{equation}

Aquí, hemos obtenido una expresión finíta para la derivada de una función evaluada en un punto, denominada \textit{derivada adelantada}. Ahora, si buscamos el error de la ecuación \eqref{eq:derad} vemos que
\begin{equation}
  Error =  \left| f^{\prime}(a) - \frac{f(a+h) - f(a)}{h} \right| = \frac{1}{2!}\left| f^{\prime \prime}(\xi) \right| h
\end{equation}
en notación $\mathcal{O}$ de Landau, 
\begin{equation}
  Error = \frac{1}{2!}\left| f^{\prime \prime}(\xi) \right| h = \mathcal{O}(h)
\end{equation}
Esto quiere decir que el error de truncamiento de la derivada adelantada será del orden de $h$, como se puede observar en la figura \ref{fig:adnterr}, el error de la derivada decrece linealmente con el orden de magnitud de $h$ hasta aproximadamente $h \sim 10^{-8}$, pasado ese punto $h$ se vuelve demasiado pequeño para la computadora y predomina el error computacional. Por esta razón, a la hora de implementar este, o cualquiera de los métodos expuestos en este apunte, es conveniente utilizar un $h \sim 10^{-5}$. 

\insertimage[\label{fig:adnterr}]{./img/derivadas/adnterr.pdf}{scale = 0.75}{Error de la derivada adelantada de $f(x) = \sin(x)$ en $x = 0.3$ para distintos valores de $h$}

\subsubsection{Otros métodos de derivación numérica}

Si bien la derivada adelantada es útil debido a lo simple de su implementación, muchas veces queremos utilizar métodos cuyo error sea de un mayor orden de $h$, es decir, menor.

Similarmente a como obtuvimos la expresión \eqref{eq:derad}, podemos obtener una derivada \textbf{retrasada}, en efecto, consideramos la serie de taylor de $f(a-h)$,
\begin{align}
  f(a - h) &= f(a) - f^{\prime}(a)h + \frac{1}{2!}f^{\prime \prime}(\xi)h^{2} \\
  f^{\prime}(a) &= \frac{f(a) - f(a-h)}{h} + \frac{1}{2!}f^{\prime \prime}(\xi)h \label{eq:deriret}
\end{align}

Si consideramos las series de taylor de $f(a+h)$ y $f(a-h)$ expandidas haste el tercer orden,
\begin{align}
  f(a+h) &= f(a) + f^{\prime}(a)h + \frac{1}{2!}f^{\prime \prime}(a)h^{2} + \frac{1}{3!}f^{\prime \prime \prime}(\xi)h^{3} \\
  f(a-h) &= f(a) - f^{\prime}(a)h + \frac{1}{2!}f^{\prime \prime}(a)h^{2} - \frac{1}{3!}f^{\prime \prime \prime}(\xi)h^{3}
\end{align}
Ahora si restamos obtenemos,
\begin{align}
  f(a+h) - f(a-h) &= 2f^{\prime}(a)h + \frac{2}{3!}f^{\prime \prime \prime}(\xi)h^{3} \\
  2f^{\prime}(a)h &= f(a+h) - f(a-h) - \frac{2}{3!}f^{\prime \prime \prime}(\xi)h^{3} \\
  f^{\prime}(a) &= \frac{f(a+h) - f(a-h)}{2h} - \frac{1}{6}f^{\prime \prime \prime}(\xi)h^{3} \label{eq:dercen}
\end{align}

Así, obtenemos una expresión para la \textbf{derivada centrada}, con error del orden de $h^{3}$.

Podemos obtener muchos otros métodos de derivación numérica por medios similares, en general el proceso consiste en jugar con series de taylor hasta poder despejar el termino $f^{\prime}(a)$. 

\subsubsection{Derivadas para datos discretos}

En general, uno no suele tratar con funciones definidas analíticamente, si no que tendremos conjuntos de datos relacionados a alguna magnitud física, si quisieramos eventualmente encontrar la derivada de la "función" que definen estos datos, debemos aprender a tratar con puntos no necesariamente equidistantes.

Similarmente a caso anterior, dado un conjunto de puntos $x_{k}$ y una función continua y diferenciable tal que $f(x_{k}) = f_{k}$, definimos la distancia entre dos puntos como $h_{k} = x_{k} - x_{k-1}$. Podemos obtener una primera aproximación simplemente tomando la pendiente entre dos puntos, esto es
\begin{equation} \label{eq:firstdisc}
  \frac{df_{k}}{dx_{k}} \approx \frac{ f_{k}-f_{k-1}}{x_{k}-x_{k-1}}
\end{equation}
Esto es analogo a la derivada atrasada encontrada en la sección anterior, y se puede demostrar que tendrá un error del orden de $h_{k}$. Notar además que si imponemos que los nodos estén igualmente espaciados, es decir, $x_{k}-x_{k-1} = h = cte.$, la ecuación \eqref{eq:firstdisc} recupera la ecuación \eqref{eq:deriret}.

Ahora, estamos incentivados a encontrar un método de derivación discreta con error de menor orden, para ello consideramos la serie de taylor de $f_{k+1}$ y $f_{k-1}$ 
\begin{align}
  f_{k+1} &= f_{k} + \fracderivat{f_{k}}{x}h_{k+1} + \frac{1}{2!}\fracnderivat{f_{k}}{x}{2}h_{k+1}^{2} + \frac{1}{3!}\fracnderivat{f(\xi_{1})}{x}{3}h_{k+1}^{3} \label{eq:kplus} \\
    f_{k-1} &= f_{k} - \fracderivat{f_{k}}{x}h_{k} + \frac{1}{2!}\fracnderivat{f_{k}}{x}{2}h_{k}^{2} - \frac{1}{3!}\fracnderivat{f(\xi_{2})}{x}{3}h_{k}^{3} \label{eq:kmin}
\end{align}

con $x_{k-1} \leq \xi_{2} \leq x_{k} \leq \xi_{1} \leq x_{k+1}$. Despejando $\fracderivat{f_{k}}{x}$ en ambas expresiones obtenemos,
\begin{align}
  \fracderivat{f_{k}}{x}h_{k+1} &= f_{k+1} - f_{k} - \frac{1}{2!}\fracnderivat{f_{k}}{x}{2}h_{k+1}^{2} - \fracnderivat{f(\xi_{1})}{x}{3}h_{k+1}^{3} \label{eq:kone} \\
    \fracderivat{f_{k}}{x}h_{k} &= f_{k} - f_{k-1} + \frac{1}{2!}\fracnderivat{f_{k}}{x}{2}h_{k}^{2} + \fracnderivat{f(\xi_{2})}{x}{3}h_{k}^{3} \label{eq:ktwo}
\end{align}
Multiplicando \eqref{eq:kone} y \eqref{eq:ktwo} por $h_{k}/h_{k+1}$ y $h_{k+1}/h_{k}$, respectivamente, 
\begin{align}
  \fracderivat{f_{k}}{x}{}h_{k} &= \frac{(f_{k+1} - f_{k})h_{k}}{h_{k+1}} - \frac{1}{2!}\fracnderivat{f_{k}}{x}{2}h_{k}h_{k+1} - \fracnderivat{f(\xi_{1})}{x}{3}h_{k+1}^{2}h_{k} \\
    \fracderivat{f_{k}}{x}{}h_{k+1} &= \frac{(f_{k} - f_{k-1})h_{k+1}}{h_{k}} + \frac{1}{2!}\fracnderivat{f_{k}}{x}{2}h_{k}h_{k+1} + \fracnderivat{f(\xi_{2})}{x}{3}h_{k}^{2}h_{k+1}
\end{align}
Sumando,
\begin{equation}
  [h_{k+1} + h_{k}]\fracderivat{f_{k}}{x} = \frac{(f_{k+1} - f_{k})h_{k}}{h_{k+1}} + \frac{(f_{k} - f_{k-1})h_{k+1}}{h_{k}} + \frac{1}{3!} \left[ f^{\prime \prime \prime}(\xi_{2})h_{k} - f^{\prime \prime \prime}(\xi_{1})h_{k+1} \right]h_{k}h_{k+1}
\end{equation}
Finalmente,

\begin{equation}
  \fracderivat{f_{k}}{x} = \frac{(f_{k+1} - f_{k})}{h_{k+1} + h_{k}}\frac{h_{k}}{h_{k+1}} + \frac{(f_{k} - f_{k-1})}{h_{k+1} + h_{k}}\frac{h_{k+1}}{h_{k}} + \frac{1}{3!} \left[ \frac{ f^{\prime \prime \prime}(\xi_{2})h_{k} - f^{\prime \prime \prime}(\xi_{1})h_{k+1}}{h_{k+1} + h_{k}} \right]h_{k}h_{k+1}
\end{equation}
Ahora, por el teorema del valor medio, podemos reemplazar el último termino entre corchetes por una función evaluada en un punto $\xi$, con $\xi_{2} \leq \xi \leq \xi_{1}$, en efecto,

\begin{align}
  \fracderivat{f_{k}}{x} &= \frac{(f_{k+1} - f_{k})}{h_{k+1} + h_{k}}\frac{h_{k}}{h_{k+1}} + \frac{(f_{k} - f_{k-1})}{h_{k+1} + h_{k}}\frac{h_{k+1}}{h_{k}} + \frac{1}{3!} \fracnderivat{f_{k}(\xi)}{x}{3} h_{k}h_{k+1} \\
  &= \frac{(f_{k+1} - f_{k})}{h_{k+1} + h_{k}}\frac{h_{k}}{h_{k+1}} + \frac{(f_{k} - f_{k-1})}{h_{k+1} + h_{k}}\frac{h_{k+1}}{h_{k}} + \mathcal{O}(h_{k}h_{k+1})
\end{align}

Así, logramos encontrar una expresión para la derivada de una función definida punto a punto con nodos no necesariamente equidistantes, con error de orden cuadratico. Puede comprobar que si imponemos que los nodos sean equidistantes ($h_{k} = h_{k+1} = h$), se recupera la expresión \eqref{eq:dercen} para la derivada centrada (\textbf{Demuestrelo!}). Más adelante estudiaremos métodos más avanzados para tratar con distribuciones de datos con interpolaciones y splines. Para terminar, cabe mencionar que la mayor desventaja de este método es que, al utilizar 3 puntos, aplicandola a una serie de puntos, la expresión numérica de la derivada tendrá dos puntos menos que la serie original. Por esta razón, en la practica se implementa una derivada adelantada para los puntos iniciales de un set de datos, y una atrasada para los finales; la derivada centrada se utiliza solo para los puntos intermedios.

